% !TeX root = ../../../progress-sem1.tex
% 4.x ปัญหาและอุปสรรคที่พบ (semester 1)
\section{ปัญหาและอุปสรรคที่พบ}
\subsection{ปัญหาและแนวทางแก้ไข}
\thaijustify{
  \begin{enumerate}
    \item \textbf{ผลลัพธ์ของ AI ไม่สม่ำเสมอ}\\
      UI ที่โมเดลสร้างในบางครั้งใช้ component ที่ไม่อยู่ใน UI kit ใช้สีไม่ตรงกับแบรนด์ หรือไม่ผ่านมาตรฐานการเข้าถึง ทีมแก้ไขโดยเพิ่มขั้นตรวจอัตโนมัติหลังสร้างทุกครั้ง หากไม่ผ่าน ระบบส่งปัญหาที่พบกลับให้โมเดลแก้ได้สูงสุดสองครั้ง และหากยังไม่ผ่านจะเริ่มสร้าง variant นั้นใหม่อีกหนึ่งครั้ง
    \item \textbf{การสร้าง UI ใช้เวลานาน}\\
      การสร้าง variant แต่ละชุดใช้เวลาหลายนาที หากรอผลในคำขอเดียว หน้าจอจะค้างและผู้ใช้ไม่รู้ว่าระบบทำอะไรอยู่ ทีมจึงย้ายงานไปทำเบื้องหลังผ่านคิว BullMQ แสดงความคืบหน้าและความคิดของโมเดลแบบ real-time ผ่าน SSE และให้ผู้ใช้ยกเลิกงานได้
    \item \textbf{การเข้าถึงบริการ AI}\\
      CPE AI Gateway เข้าถึงได้ผ่านเครือข่าย VPN ของมหาวิทยาลัยเท่านั้น ระหว่างพัฒนาสมาชิกทุกคนจึงต้องเชื่อมต่อ VPN และเซิร์ฟเวอร์ที่ใช้ deploy ต้องเข้าถึง gateway ได้เช่นกัน
    \item \textbf{การสมัครสมาชิกด้วยอีเมลล่าช้ากว่าแผน}\\
      ทีมพัฒนาการเข้าสู่ระบบด้วยบัญชี Google เสร็จก่อน ส่วนการสมัครด้วยอีเมลและรหัสผ่าน การยืนยันอีเมล และการรีเซ็ตรหัสผ่าน อยู่ระหว่างพัฒนาใน branch แยก และจะรวมเข้า branch หลักภายในภาคการศึกษาที่ 1
  \end{enumerate}
  \todo{ปัญหาด้านการทำงานเป็นทีม เวลา หรือทรัพยากร (ถ้ามี) และยืนยันรายการข้างต้น}
}

\subsection{การเปลี่ยนแปลงจากข้อเสนอโครงงาน}
\thaijustify{
  ระหว่างออกแบบหน้าจอและพัฒนาระบบ ทีมปรับรายละเอียดบางส่วนจากข้อเสนอโครงงาน ดังนี้
  \begin{enumerate}
    \item \textbf{เปลี่ยน organization เป็น workspace}\\
      ผู้ใช้สร้างและสลับ workspace ได้หลายแห่ง และหนึ่ง workspace มีหลายโปรเจกต์
    \item \textbf{ผู้เข้าร่วมทดสอบไม่ต้องมีบัญชี}\\
      ผู้เข้าร่วมทดสอบเข้าเซสชันด้วยรหัสหรือลิงก์และกรอกเพียงชื่อและอีเมล โดยรหัสหนึ่งใช้ได้กับผู้เข้าร่วมทดสอบหนึ่งคน
    \item \textbf{กำหนดอุปกรณ์ที่ทดสอบต่อ user story}\\
      เดิมกำหนด platform ไว้ที่ระดับโปรเจกต์ แต่ user story แต่ละเรื่องอาจต้องการทดสอบบนอุปกรณ์ต่างกัน จึงย้ายไปเลือกในหน้า user story
    \item \textbf{ปรับรอบการทดสอบตามผู้เข้าร่วมทดสอบแต่ละคน}\\
      รอบแรกผู้เข้าร่วมทดสอบทุกคนทดสอบ A/B ชุดเดียวกัน ตั้งแต่รอบที่สอง AI สร้าง A/B ของรอบถัดไปจากข้อเสนอแนะของผู้เข้าร่วมทดสอบคนนั้น โดยนักออกแบบเข้าไปกำหนดทิศทางเมื่อแนวโน้มความคิดเห็นขึ้นลงผิดปกติ และตรวจผลรวมเมื่อทุกคนทดสอบเสร็จ เพื่อลดภาระของนักออกแบบในระหว่างการทดสอบ
    \item \textbf{ปรับการนำทางของหน้าจอ}\\
      เปลี่ยนจากแถบนำทางด้านบนและตัวบอกขั้นตอน (stepper) เป็น sidebar ที่แสดงขั้นตอนของ user story เนื่องจากแถบด้านบนซ้ำซ้อนกับเมนูอื่น และ stepper ทำให้ผู้ใช้สับสนเมื่อเริ่มจากขั้นกลาง
  \end{enumerate}
  \todo{เพิ่มเหตุผลของการเปลี่ยนแปลงข้อ 1 และ 2 และความเห็นกรรมการที่เกี่ยวข้อง (ถ้ามี)}
}
